\documentclass[10pt]{article}
\usepackage[margin=0.72in]{geometry}
\usepackage{amsmath}
\usepackage{amssymb}
\usepackage{booktabs}
\usepackage{graphicx}
\usepackage{caption}
\usepackage{microtype}
\title{SAC versus action-conditioned JEPA on the fractal terrain task}
\author{}
\date{RTX 3090 runs, September 20, 2026}
\begin{document}
\maketitle

\section*{Question and result}
All runs use the same centered-Gaussian squared-potential terrain, random
starts, $P_{\max}=6$ longitudinal power cap, 16 environments, held-out terrain
seeds, and 10,000-transition budget. SAC alone is the strongest observed
controller. The normalized, action-conditioned JEPA branch is mechanically
valid---it beats shuffled-action and identity diagnostics---but still reduces
held-out speed at this budget.

\begin{center}
\small
\begin{tabular}{p{1.6in}ccp{2.0in}}
\toprule
Run & Held-out speed at 10k & Difference from SAC & Status \\
\midrule
SAC-only & $1.655\pm0.056$ & --- & baseline \\
SAC+JEPA, old MSE target & $1.278\pm0.042$ & $-22.8\%$ & superseded cadence/scale safeguards \\
SAC+JEPA, corrected MSE & $1.403\pm0.020$ & $-15.2\%$ & variance floor and slow EMA \\
SAC+JEPA, cosine predictive loss & $1.188$ & $-28.2\%$ & scale-normalized diagnostic run \\
\bottomrule
\end{tabular}
\end{center}

\section*{Full corrected SAC+JEPA loss}
For a replayed transition $(o_t,a_t,o_{t+1})$, the actor has online encoder
$z_t=f_\theta(o_t)$, EMA target $\bar z_{t+1}=\bar f(o_{t+1})$, and predictor
$q_\psi$. The optimized actor objective is
\begin{equation}
\begin{aligned}
\mathcal L_{\rm actor}={}&\mathbb E\!\left[\alpha\log\pi(a|o)-\min(Q_1(o,a),Q_2(o,a))\right]\\
&+\lambda\left(1-\cos\!\left(q_\psi([z_t,a_t]),\operatorname{sg}[\bar z_{t+1}]\right)\right)\\
&+c_{\rm var}\frac{1}{d}\sum_{j=1}^d\max\!\left(0,s_{\rm min}-\operatorname{std}_B(z_{:,j})\right).
\end{aligned}
\end{equation}
The cosine loss is scale-free; the variance floor prevents a zero-variance
representation. The final run uses $\lambda=0.025$, $c_{\rm var}=0.01$,
$s_{\rm min}=0.1$, and updates the EMA target by 0.01 every 64 collected
transitions. Logged diagnostics include the weighted JEPA and variance terms,
action-shuffled loss, identity baseline, and embedding spread.

\begin{figure}[p]
\centering
\includegraphics[width=.49\linewidth]{../output/pdf/sac_only_learning_curve.png}
\includegraphics[width=.49\linewidth]{../dumps/sac_jepa_cosine_10k/plots/learning_curve.png}
\caption{Performance curves. Left: SAC-only. Right: normalized cosine SAC+JEPA.
The right lower panel is the scale-free predictive loss.}
\end{figure}

\begin{figure}[p]
\centering
\includegraphics[width=.49\linewidth]{../dumps/sac_10k/plots/latent_umap_sac_000010000_seed100001.png}
\includegraphics[width=.49\linewidth]{../dumps/sac_jepa_cosine_10k/plots/latent_umap_sac_000010000_seed100001.png}
\caption{Endpoint terrain-overlay trajectories and UMAPs on the same held-out
seed. Left: SAC-only. Right: normalized cosine SAC+JEPA.}
\end{figure}

\begin{figure}[p]
\centering
\includegraphics[width=.49\linewidth]{../dumps/sac_10k/plots/dynamics_history_sac_000010000_seed100001.png}
\includegraphics[width=.49\linewidth]{../dumps/sac_jepa_cosine_10k/plots/dynamics_history_sac_000010000_seed100001.png}
\caption{Full velocity, control, and acceleration diagnostics. Left: SAC-only.
Right: normalized cosine SAC+JEPA.}
\end{figure}

\begin{figure}[p]
\centering
\includegraphics[width=.49\linewidth]{../dumps/sac_10k/plots/encoder_cka_matrix_seed100001.png}
\includegraphics[width=.49\linewidth]{../dumps/sac_jepa_cosine_10k/plots/encoder_cka_matrix_seed100001.png}
\caption{Encoder evolution across the five saved checkpoints, measured by
linear CKA on a common held-out inference trajectory.}
\end{figure}

\begin{figure}[p]
\centering
\includegraphics[width=.49\linewidth]{../dumps/sac_10k/plots/pairwise_encoder_umap_5x5_seed100001.png}
\includegraphics[width=.49\linewidth]{../dumps/sac_jepa_cosine_10k/plots/pairwise_encoder_umap_5x5_seed100001.png}
\caption{Checkpoint-pair latent matrix on the shared deterministic held-out
trajectory. Each cell jointly embeds the row and counterpart encoder states
(blue/orange); diagonal cells are self-comparisons. Left: SAC-only. Right:
normalized cosine SAC+JEPA.}
\end{figure}

\section*{Interpretation}
The cosine replacement fixes the earlier loss-scale failure: the predictive
term is no longer numerically negligible relative to the variance term, and
the shuffled-action loss exceeds the real-action loss ($0.0099$ versus $0.0038$
at 10k). Thus the predictor uses actions. Nonetheless, performance is lower
than SAC alone. At this horizon and budget, the evidence favors SAC-only;
multi-step targets and matched multi-seed sweeps remain the next JEPA test.

\clearpage
\section*{V2: five-times-longer SAC+JEPA run (50k transitions)}
V2 repeats the corrected cosine SAC+JEPA configuration from V1 on the RTX
3090, with the same terrain split, power limit, random starts, and held-out
evaluation seeds. Only the transition budget changes: 10k to 50k. Checkpoints
are saved every 10k, so all representation plots compare five within-run
checkpoints on the same held-out deterministic trajectory.

\begin{center}
\small
\begin{tabular}{lcc}
\toprule
Run & Held-out speed & Interpretation \\
\midrule
V1 cosine SAC+JEPA, 10k & $1.188$ & early-budget endpoint \\
V2 cosine SAC+JEPA, 50k & $2.954\pm0.049$ & continued improvement, plateau after about 18k \\
Matched SAC-only, 50k & running & required before asserting relative slowness \\
\bottomrule
\end{tabular}
\end{center}

The auxiliary branch remains active in V2: at 10,240 transitions its real
action-conditioned cosine loss is $0.0045$ versus $0.0101$ with actions
shuffled. The gap narrows late in training, as the controller and observations
become more predictable. Thus V2 establishes that the JEPA controller keeps
learning with a larger budget, but it does \emph{not} yet prove that it is
slower than SAC: that requires the matched 50k SAC-only control now in flight.

\begin{figure}[p]
\centering
\includegraphics[width=.92\linewidth]{../dumps/sac_jepa_cosine_50k_v2/plots/learning_curve.png}
\caption{V2 held-out performance and auxiliary diagnostics across 50k
transitions. The held-out speed rises sharply after 10k and stabilizes near
2.95 from approximately 18k onward.}
\end{figure}

\begin{figure}[p]
\centering
\includegraphics[width=.49\linewidth]{../dumps/sac_jepa_cosine_50k_v2/plots/latent_umap_sac_000050000_seed100001.png}
\includegraphics[width=.49\linewidth]{../dumps/sac_jepa_cosine_50k_v2/plots/dynamics_history_sac_000050000_seed100001.png}
\caption{V2 endpoint deterministic inference on held-out seed 100001. Left:
terrain-overlay trajectory and encoder UMAP. Right: velocity, input, field,
damping, and net acceleration histories.}
\end{figure}

\begin{figure}[p]
\centering
\includegraphics[width=.49\linewidth]{../dumps/sac_jepa_cosine_50k_v2/plots/encoder_cka_matrix_seed100001.png}
\includegraphics[width=.49\linewidth]{../dumps/sac_jepa_cosine_50k_v2/plots/pairwise_encoder_umap_5x5_seed100001.png}
\caption{V2 encoder evolution across the five 10k-spaced checkpoints on one
shared held-out trajectory. Left: linear CKA. Right: pairwise joint UMAP
matrix (blue/orange are the two checkpoint encoders in each cell).}
\end{figure}

\clearpage
\section*{V2 completion: matched 50k SAC-only control}
The control mentioned above has now completed under the same 50k transition
budget, terrain split, power limit, random-start protocol, evaluation seeds,
and checkpoint schedule. This is the valid comparison for the question
``does JEPA merely learn more slowly?''

\begin{center}
\small
\begin{tabular}{lccc}
\toprule
Transitions & SAC-only & SAC+JEPA & JEPA minus SAC \\
\midrule
10,240 & $1.426$ & $1.402$ & $-0.024$ \\
12,288 & $2.258$ & $2.433$ & $+0.175$ \\
16,384 & $2.836$ & $2.844$ & $+0.008$ \\
50,000 & $2.952\pm0.030$ & $2.954\pm0.049$ & $+0.003$ \\
\bottomrule
\end{tabular}
\end{center}

\begin{figure}[h]
\centering
\includegraphics[width=.88\linewidth]{../dumps/sac_jepa_cosine_50k_v2/plots/matched_sac_vs_jepa_50k.png}
\caption{Direct matched 50k held-out comparison. The curves cross during the
early transient and are nearly identical from about 16k onward. Shading is the
held-out standard deviation across the fixed evaluation seeds.}
\end{figure}

\noindent\textbf{Conclusion.} In this one matched seed/configuration pair,
the corrected cosine SAC+JEPA branch is neither persistently worse nor
detectably slower than SAC-only. Its final mean differs by $0.003$ speed
($0.09\%$), far smaller than the reported evaluation spread. This establishes
convergence at 50k; it does not establish a general JEPA benefit. Matched
multi-seed replicates are required for that claim.

\clearpage
\section*{V2 visual companion: SAC-only versus SAC+JEPA}
These panels are the corresponding standard inference artifacts from the two
matched 50k endpoints. Both use the same held-out terrain seed (100001) and
the same deterministic-inference protocol; only the learned controller and
encoder differ.

\begin{figure}[p]
\centering
\includegraphics[width=.49\linewidth]{../dumps/sac_50k_control_v2/plots/latent_umap_sac_000050000_seed100001.png}
\includegraphics[width=.49\linewidth]{../dumps/sac_jepa_cosine_50k_v2/plots/latent_umap_sac_000050000_seed100001.png}
\caption{Matched endpoint terrain-overlay trajectory and latent UMAP. Left:
SAC-only. Right: SAC+JEPA. Point colour maps the two UMAP coordinates back to
the physical trajectory.}
\end{figure}

\begin{figure}[p]
\centering
\includegraphics[width=.49\linewidth]{../dumps/sac_50k_control_v2/plots/dynamics_history_sac_000050000_seed100001.png}
\includegraphics[width=.49\linewidth]{../dumps/sac_jepa_cosine_50k_v2/plots/dynamics_history_sac_000050000_seed100001.png}
\caption{Matched endpoint velocity and acceleration histories. Left:
SAC-only. Right: SAC+JEPA. Each panel contains speed and latent velocity,
actions, and signed input/potential/damping/net acceleration components.}
\end{figure}

\begin{figure}[p]
\centering
\includegraphics[width=.49\linewidth]{../dumps/sac_50k_control_v2/plots/encoder_cka_matrix_seed100001.png}
\includegraphics[width=.49\linewidth]{../dumps/sac_jepa_cosine_50k_v2/plots/encoder_cka_matrix_seed100001.png}
\caption{Within-run representation evolution on a common held-out trajectory.
Left: SAC-only. Right: SAC+JEPA. Each 5 by 5 matrix is linear CKA across the
10k, 20k, 30k, 40k, and 50k checkpoints.}
\end{figure}

\begin{figure}[p]
\centering
\includegraphics[width=.49\linewidth]{../dumps/sac_50k_control_v2/plots/pairwise_encoder_umap_5x5_seed100001.png}
\includegraphics[width=.49\linewidth]{../dumps/sac_jepa_cosine_50k_v2/plots/pairwise_encoder_umap_5x5_seed100001.png}
\caption{Pairwise joint UMAP trajectory matrices for the same checkpoints.
Blue and orange are the row and counterpart encoders evaluated on identical
held-out inference states. Left: SAC-only. Right: SAC+JEPA.}
\end{figure}

\clearpage
\section*{Post-run diagnostic fidelity correction}
The original V2 run is unchanged. Its speed curve, endpoint evaluations, and
inference/representation artifacts remain valid. Two reported auxiliary
details are corrected for future runs: (i) the cosine-prediction panel is now
labelled ``cosine distance ($1-\cos$)'', not MSE; and (ii) training now logs
the mean of every replay update since the prior evaluation, rather than the
final minibatch alone. Consequently, the V2 single-checkpoint real-action
versus shuffled-action loss values are diagnostic snapshots, not evidence of
a statistically significant action effect. The corrected identity baseline
also compares target-encoder embeddings of $o_t$ and $o_{t+1}$, isolating
temporal predictability from EMA lag. These corrections do not change the
matched 50k performance conclusion; they improve fidelity for subsequent
JEPA ablations.

\end{document}
